\documentclass[aspectratio=169,xcolor=dvipsnames]{beamer}

\usepackage{amsthm}
\usepackage{amsmath}
\usepackage{amssymb}
\usepackage{tikz}
\usepackage{pgfplots}
\usepackage{xfp}
\pgfplotsset{compat=1.18}
\usepackage[authoryear, round]{natbib}
\usepackage{hyperref}
\usetikzlibrary{arrows.meta,positioning,decorations.pathreplacing}
\hypersetup{
    colorlinks=true,
    linkcolor=blue,
    citecolor=blue,
    urlcolor=blue
}

\definecolor{darkgreen}{rgb}{0,0.5,0}
\definecolor{darkblue}{rgb}{0,0,0.55}

\newtheorem{proposition}{Proposition}
\setbeamertemplate{footline}[frame number]
\usetheme{Frankfurt}
\usefonttheme{serif}

% Title info
\title{ECON8037: Financial Economics}
\subtitle{\normalsize Week 7 Lecture - The Economics of Cryptocurrencies I: Blockchain, Mining, and Double-Spending (Chiu \& Koeppl, 2017)}
\author{Simon Mishricky}
\institute{Research School of Economics, Australian National University}
\date{\today}

\begin{document}
\section{Motivation}
\frame{\titlepage}

\begin{frame}{Why study the economics of cryptocurrencies?}
    Cryptocurrencies (Bitcoin, Ethereum, $\ldots$) are an attempt to support payments
    \textbf{without} a trusted third party.
    \vspace{2mm}
    \begin{itemize}
        \item Critics: fraud, bubbles, vehicle for illegal activity.
        \item Advocates: cryptographically secure decentralised payments.
        \item Central banks: actively exploring blockchain-based digital currencies.
    \end{itemize}
    \vspace{2mm}
    \citet{chiu2017} ask: \emph{how well can a cryptocurrency serve as a means of payment?}
    \vspace{2mm}
    \begin{itemize}
        \item Build a general equilibrium monetary model with a blockchain;
        \item discipline it with Bitcoin transactions data;
        \item compute the welfare cost of using Bitcoin and an optimally designed alternative.
    \end{itemize}
\end{frame}

\begin{frame}{Roadmap for Lecture I}
    \begin{enumerate}
        \item Cryptocurrencies in one slide: blockchain, mining, PoW.
        \item The double-spending problem.
        \item A partial-equilibrium model of mining and secret mining.
        \item Lagos--Wright general equilibrium with a cryptocurrency.
        \item Double-spending-proof equilibrium.
    \end{enumerate}
    \vspace{4mm}
    Lecture II picks up with optimal design, welfare, and applications.
\end{frame}

\begin{frame}{Three key takeaways}
    \begin{enumerate}
        \item A cryptocurrency must overcome \textbf{double spending} by combining \emph{costly mining}
              (proof-of-work) with \emph{settlement lags} (confirmations).
        \vspace{2mm}
        \item This creates a fundamental trade-off: settlement cannot be both
              \textbf{immediate} and \textbf{final}.
        \vspace{2mm}
        \item The model embeds these features in a Lagos--Wright environment so that
              the value of the currency, mining effort, and double-spending incentives
              are jointly determined.
    \end{enumerate}
\end{frame}

\section{Cryptocurrencies overview}

\begin{frame}{Digital currency vs.\ cryptocurrency}
    A \textbf{digital token} is a string of bits. Digital strings can be copied costlessly
    $\Rightarrow$ the \emph{double-spending problem}.
    \vspace{3mm}
    \begin{itemize}
        \item \textbf{Centralised solution (e.g.\ PayPal):} a trusted third party
        maintains a ledger and updates it as people transact. Users must trust the intermediary.
        \vspace{2mm}
        \item \textbf{Decentralised solution (e.g.\ Bitcoin):} a network of anonymous validators
        maintains and updates copies of the ledger. Trust is replaced by a \emph{consensus protocol}.
    \end{itemize}
    \vspace{3mm}
    The cryptocurrency must induce validators (i) to keep an accurate ledger and (ii) to
    deter users from double spending.
\end{frame}

\begin{frame}{Blockchain as a record-keeping device}
    A \textbf{block} = a set of transactions conducted between users in a period.
    A \textbf{blockchain} = a chain of blocks recording all past transactions.
    \vspace{3mm}
    \begin{itemize}
        \item Each block is linked to the previous one (dynamically consistent ledger).
        \item Anyone can verify any user's balance by reading the chain.
        \item Different from traditional money, which only records the \emph{current}
              distribution of balances.
    \end{itemize}
    \vspace{3mm}
    Consensus rule: the \textbf{longest chain} is the trusted public record.
\end{frame}

\begin{frame}{Mining and Proof-of-Work}
    \textbf{Mining} = competition to update the blockchain with a new block.
    \vspace{2mm}
    \begin{itemize}
        \item A miner who first solves a computationally costly puzzle (\emph{Proof-of-Work, PoW})
        wins the right to append a new block.
        \item Winning probability is proportional to the miner's fraction of computational power.
        \item Winners receive a reward $R$, financed by (i) newly created coins (seigniorage)
              and (ii) transaction fees.
    \end{itemize}
    \vspace{3mm}
    \textbf{Why is mining costly?} To make rewriting history expensive. Older transactions
    are more trustworthy because revising them requires re-doing the PoW for every subsequent block.
\end{frame}

\begin{frame}{The bootstrapping problem}
    A cryptocurrency must jointly support three things:
    \vspace{2mm}
    \begin{itemize}
        \item \textbf{Security of the blockchain} -- so merchants accept the currency;
        \item \textbf{Value of the currency} -- so mining rewards are valuable;
        \item \textbf{Mining activity} -- so double spending is costly.
    \end{itemize}
    \vspace{3mm}
    Each side reinforces the others. The model must determine all three jointly in equilibrium.
\end{frame}

\section{Double Spending}

\begin{frame}{The double-spending attack: forward looking}
    Even with a blockchain, a \emph{forward-looking} attack is possible:
    \vspace{2mm}
    \begin{enumerate}
        \item Buyer sends payment $d$ to seller; instruction is broadcast to miners.
        \item Buyer secretly mines an \emph{alternative} block in which the payment does not occur.
        \item Once seller delivers the goods, buyer releases the secret chain.
        \item If the secret chain is longer, the consensus protocol adopts it and the
              original payment is wiped out.
    \end{enumerate}
    \vspace{2mm}
    Outcome: buyer keeps both the goods and the balances; seller is left empty handed.
\end{frame}

\begin{frame}{Confirmation lags as a defence}
    Defence: seller waits for $N$ \emph{confirmations} (i.e.\ $N$ blocks added on top of the
    block containing the payment) before delivering goods.
    \vspace{3mm}
    \begin{itemize}
        \item To double spend, buyer must mine \textbf{$N+1$ blocks in a row} faster than the
              rest of the network.
        \item Longer $N$ $\Rightarrow$ harder to attack, but slower settlement.
    \end{itemize}
    \vspace{3mm}
    Trust = (security of blockchain) $\times$ (mining ecosystem) $\times$ (value of currency).
\end{frame}

\section{Model setup}

\begin{frame}{Partial equilibrium: one transaction period}
    Within a single period there are $\bar N+1$ subperiods.
    Buyer holds real balances $z$; bargains with a seller over $(x,d,N)$:
    \vspace{2mm}
    \begin{itemize}
        \item $x$ = quantity of goods,
        \item $d\le z$ = payment in real balances,
        \item $N\in\{0,\ldots,\bar N\}$ = confirmation lag.
    \end{itemize}
    \vspace{3mm}
    Buyer's preferences (with confirmation lag $N$):
    \[
        \delta^N u(x),\qquad \delta\in(0,1).
    \]
    Across periods agents discount by $\beta = \delta^{\bar N+1}$.
\end{frame}

\begin{frame}{The mining game}
    $M$ miners each invest computing power $q(i)$ measured in real balances.
    A particular miner wins block $n$ with probability
    \[
        \rho(i) \;=\; \frac{q(i)}{\sum_{m=1}^M q(m)}. \tag{3}
    \]
    Each subperiod, each miner solves
    \[
        \max_{q(i)} \;\; \rho_i\,\beta R - q(i). \tag{4}
    \]
    Symmetric Nash equilibrium $q(m)=Q$ for all $m$:
    \[
        Q \;=\; \frac{M-1}{M^2}\beta R. \tag{6}
    \]
\end{frame}

\begin{frame}{Free entry of miners}
    Total computing cost per subperiod:
    \[
        MQ \;=\; \frac{M-1}{M}\beta R. \tag{7}
    \]
    Take $M\to\infty$ (competitive mining):
    \vspace{2mm}
    \begin{lemma}[Rent dissipation]
        As $M\to\infty$, miners earn zero profits and aggregate computing power dissipates all rewards:
        \[
            MQ \;=\; \beta R.
        \]
    \end{lemma}
    \vspace{2mm}
    Mining is a pure rent-seeking activity: every dollar of reward is matched by a dollar of cost.
\end{frame}

\begin{frame}{Secret mining and double-spending payoffs}
    To double spend successfully, the buyer must win the mining game in $N+1$ consecutive subperiods.
    Conditional on $N$ prior successes, his subperiod-$N$ payoff from investing $q_N$ is
    \[
        \rho(q_N)\,\beta[d + (N+1)R] - q_N. \tag{9}
    \]
    Define
    \[
        \Delta \;\equiv\; \frac{d}{R} + (N+1). \tag{11}
    \]
    As $M\to\infty$:
    \[
        \hat q_N(d,N) = \beta R\bigl[\sqrt{\Delta} - 1\bigr], \qquad
        \rho(\hat q_N) = \frac{\sqrt{\Delta}-1}{\sqrt{\Delta}}. \tag{12,13}
    \]
\end{frame}

\begin{frame}{Backward induction}
    Define recursively the expected double-spending payoff in subperiod $n$:
    \[
        D_n(d,N) \;=\; \max_{q_n} \rho(q_n)\, D_{n+1}(d,N) - q_n, \qquad
        D_N(d,N) = \beta R(\sqrt\Delta-1)^2.
    \]
    \begin{lemma}[Closed form]
        For $s=0,\ldots,N$,
        \[
            D_{N-s}(d,N) = \beta R[\sqrt\Delta - (s+1)]^2,\quad
            \hat q_{N-s} = \beta R[\sqrt\Delta - (s+1)].
        \]
    \end{lemma}
    \vspace{2mm}
    Investment is increasing in $n$: if it's optimal to engage in secret mining today,
    it's optimal to continue tomorrow conditional on success.
\end{frame}

\begin{frame}{The no-double-spending (DS) constraint}
    Double spending is unattractive whenever
    \[
        \beta R\bigl[\sqrt\Delta - (N+1)\bigr] < 0,
    \]
    so $\hat q_0=0$ and $D_0=0$. This yields the central inequality.
    \vspace{3mm}
    \begin{proposition}[DS-proof contract]
        As $M\to\infty$, a contract $(d,N)$ is double-spending proof iff
        \[
            \boxed{\;d \;<\; R\,(N+1)\,N.\;} \tag{20}
        \]
    \end{proposition}
    \vspace{2mm}
    Intuition:
    \begin{itemize}
        \item Larger payment $d$ $\Rightarrow$ stronger incentive to attack.
        \item Larger reward $R$ or longer lag $N$ $\Rightarrow$ harder to attack.
    \end{itemize}
\end{frame}

\begin{frame}{Probability of successful double spending}
    More generally, whenever (20) fails, the unconditional probability of a successful attack is
    \[
        P(d,N) \;=\; \frac{\sqrt\Delta - (N+1)}{\sqrt\Delta},\qquad \Delta = \frac{d}{R} + (N+1).
    \]
    \vspace{2mm}
    \textbf{Settlement concepts.}
    \begin{itemize}
        \item Settlement is \emph{immediate} if $N=0$, \emph{delayed} if $N>0$;
        \item Settlement is \emph{final} if $P(d,N)=0$, \emph{probabilistic} if $P>0$.
    \end{itemize}
\end{frame}

\begin{frame}{Fundamental impossibility}
    \begin{theorem}[Chiu--Koeppl]
        For any cryptocurrency based on a PoW protocol, settlement cannot be both
        \emph{immediate} and \emph{final}.
    \end{theorem}
    \vspace{3mm}
    Why? At $N=0$ the DS constraint $d<R(N+1)N$ fails for any $d>0$.
    Final settlement requires $N\ge 1$.
    \vspace{3mm}
    This is a built-in cost of decentralised consensus: trust must be earned by waiting.
\end{frame}

\section{General equilibrium}

\begin{frame}{Embedding the model in Lagos--Wright}
    \citet{chiu2017} embed the partial equilibrium model in a standard
    Lagos--Wright (2005) environment.
    \vspace{2mm}
    \begin{itemize}
        \item Time $t=0,1,\ldots$; continuum of buyers $B$, sellers $S=B\sigma$, $\sigma\in(0,1)$.
        \item Each period: \textbf{day market} (centralised, general good $h$) then
              \textbf{night market} ($\bar N+1$ subperiods).
        \item In the first subperiod, buyers match with sellers w.p.\ $\sigma$.
        \item Buyer's preference shock $\varepsilon\sim F_\varepsilon$ at start of period.
        \item Utility from night-market consumption: $\delta^N \varepsilon\, u(x)$.
    \end{itemize}
\end{frame}

\begin{frame}{Cryptocurrency in general equilibrium}
    The cryptocurrency $m$ is a digital record of balances.
    \vspace{2mm}
    \begin{itemize}
        \item Day market: perfect monitoring, no double spending.
        \item Night market: anonymous; cryptocurrency is the only means of payment.
        \item PoW protocol updates the chain each subperiod.
    \end{itemize}
    \vspace{2mm}
    The mining reward is endogenous and pinned down by design parameters:
    \[
        R \;=\; \frac{Z(\mu - 1) + D\tau}{\bar N+1}, \tag{23}
    \]
    where $\mu\ge 1$ is the gross money growth rate, $\tau\ge 0$ the transaction fee,
    $Z=BE(z)$ aggregate balances, $D=B\sigma E(d)$ aggregate transfers.
\end{frame}

\begin{frame}{Day market value function}
    Let $z=\phi m$ be real balances. The day-market problem of a buyer is
    \[
        w(z) \;=\; \max_{z',h} -h + v(z',\varepsilon)\quad
        \text{s.t.}\quad z+h \ge z' \ge 0. \tag{24,25}
    \]
    Linearity of utility in $h$ implies
    \[
        w(z) \;=\; z + w(0),\qquad
        \mathbb{E}[w(z,\varepsilon)] = z + W. \tag{27,28}
    \]
    Standard Lagos--Wright result: balances enter linearly into the value function.
\end{frame}

\begin{frame}{Night market and buyer's problem}
    A matched buyer holding balances $z$ makes a TIOLI offer $(x,d,N)$ to the seller.
    Buyer's night-market value:
    \[
        v(z;\varepsilon) = \frac{\beta}{\mu}z + \sigma[\delta^N\varepsilon u(x) + D_0(d,N) - \tfrac{\beta}{\mu}d] + \beta W. \tag{29}
    \]
    Seller's value:
    \[
        v(z;0) = \tfrac{\beta}{\mu}d(1-\tau)(1-P(d,N)) - x + \tfrac{\beta}{\mu}z + \beta W. \tag{30}
    \]
    \vspace{1mm}
    Cash-in-advance: with $\mu\ge 1 > \beta$, sellers carry no balances and buyers carry exactly $z'=d$.
\end{frame}

\begin{frame}{Buyer's optimal contract}
    Buyer solves
    \[
        \max_{(x,d,N)} \;-d + (1-\sigma)\frac{\beta}{\mu}d + \sigma\bigl[\delta^N\varepsilon u(x) + D_0(d,N)\bigr]
    \]
    subject to the seller's participation constraint:
    \[
        x \;=\; \frac{\beta}{\mu}d\,(1-\tau)\,(1-P(d,N)). \tag{33}
    \]
    \vspace{3mm}
    Two possible regimes:
    \begin{itemize}
        \item DS-proof: $d\le R(N+1)N$, so $P=D_0=0$.
        \item DS-attack: $P>0$, $D_0>0$ -- attractive only when balances are very costly to carry.
    \end{itemize}
\end{frame}

\begin{frame}{DS-proof equilibrium}
    Aggregate transfers and money demand:
    \[
        D = B\sigma\!\int_0^\infty z^*(\varepsilon;R)\,dF_\varepsilon(\varepsilon),\quad
        Z = B\!\int_0^\infty z^*(\varepsilon;R)\,dF_\varepsilon(\varepsilon). \tag{35,36}
    \]
    \vspace{2mm}
    \begin{definition}[DS-proof cryptocurrency equilibrium]
        Given $(\mu,\tau)$: offers $(x(\varepsilon),d(\varepsilon),N(\varepsilon))$, money demand $z(\varepsilon)>0$ and mining choice $q$ such that
        \begin{enumerate}
            \item the buyer's contract maximises utility;
            \item mining is a Nash equilibrium in every subperiod;
            \item the contract satisfies $d \le R(N+1)N$;
            \item the day market for balances clears.
        \end{enumerate}
    \end{definition}
\end{frame}

\begin{frame}{Existence}
    \begin{proposition}[Existence]
        If $B$ is sufficiently high, a DS-proof cryptocurrency equilibrium exists.
    \end{proposition}
    \vspace{3mm}
    The key insight:
    \begin{itemize}
        \item Mining reward $R$ is financed by \textbf{aggregate} transactions
              (scales with $B$).
        \item Incentive to double spend depends only on \textbf{individual} transaction
              size $d$.
        \item So large $B$ $\Rightarrow$ large $R$ $\Rightarrow$ DS constraint $d<R(N+1)N$
              is easy to satisfy.
    \end{itemize}
    \vspace{3mm}
    Cryptocurrencies are essentially \emph{scale-dependent}: they work better when the
    user base is large.
\end{frame}

\begin{frame}{Summary of Lecture I}
    \begin{enumerate}
        \item Cryptocurrency = digital token + blockchain + PoW + consensus rule.
        \item Double-spending attacks deterred by (i) costly mining and (ii) confirmation lags.
        \item Closed-form attack technology gives a tractable DS constraint $d < R(N+1)N$.
        \item Free-entry of miners $\Rightarrow$ all rewards dissipated as mining cost ($MQ=\beta R$).
        \item Lagos--Wright embedding determines $R$ jointly with the value of the currency
              and individual contracts.
        \item Settlement is never simultaneously immediate \emph{and} final.
    \end{enumerate}
    \vspace{4mm}
    Next lecture: optimal design ($\mu,\tau$), calibration to Bitcoin, welfare, applications.
\end{frame}

\begin{frame}[allowframebreaks]{References}
    \small
    \begin{thebibliography}{99}
        \bibitem[Chiu and Koeppl(2017)]{chiu2017}
            Chiu, J.\ and T.\ V.\ Koeppl (2017). The Economics of Cryptocurrencies --
            Bitcoin and Beyond. \textit{Bank of Canada Staff Working Paper / Queen's University}.
        \bibitem[Lagos and Wright(2005)]{lagos2005}
            Lagos, R.\ and R.\ Wright (2005). A Unified Framework for Monetary Theory
            and Policy Analysis. \textit{Journal of Political Economy} 113(3), 463--484.
    \end{thebibliography}
\end{frame}

\end{document}
